\ssscp{Distribution of glosses of PMP *duRi ‘thorn’ > ‘bone’}{15-04-02-03}
PMP *duRi is reconstructed meaning ‘thorn’.
There are a few languages west
and north of the Blust line whose reflexes are glossed ‘fish
bone’ or ‘bone’, but ‘thorn’ is clearly the dominant gloss.\footnote{
		For example,
		\citet{BlustTrussel2020} list ‘bone’ as the primary gloss for
		Philippine and North Sulawesi languages \ili{Kalamian Tagbanwa},
		Agutaynen, \ili{Sangil}, \ili{Tonsea}, and \ili{Tonsawang} west of the Blust line.}
East of the Blust line the gloss shifts to the general word
for ‘bone’ (it is the word given when talking about human body
parts, or pig anatomy), and this continues out into Oceania.
See \frf{15-02}.

\begin{figure}[h]\centering
	\includegraphics[width=\columnwidth]{duRi.pdf}
	\fcp{Distribution of glosses of PMP *duRi}{15-02}
\end{figure}

\begin{figure}[h]\centering
	\includegraphics[width=0.3\columnwidth]{15-03.jpg}
	\fcp{\ili{Buru} \it{lakat} ‘thorn’}{15-03}
\end{figure}

But a few languages in our target region seem to also retain
the gloss for ‘thorn’ alongside ‘bone’.
And some languages, such as \ili{Bima} (\it{rui} ‘thorn’,
\it{peke} ‘bone’), \ili{Helong} (\it{duliʔ} ‘thorn’, \it{seen} ‘bone’),
and \tsc{Aru} languages (Proto-\ili{Aru} **ruRi ‘thorn’ [\trf{06-04}],
**tul- ‘bone’) retain the gloss of ‘thorn’
and show no sign of the meaning ‘bone’.
Are the \ili{Helong} and \tsc{Aru} glosses a retention of the older
meaning of ‘thorn’? Or superstrate relexification from \ili{Malay}
\it{duri} ‘thorn’? Or a bit of both? The consistency of sound
correspondences favours retention, but relexification can add gentle pressure on top of that.
And once again, western
and central \tsc{\ili{Bima}-Lembata} does not regularly participate
in this widespread shift.

Many languages \ili{east} of the Blust line have completely different
words for thorn (e.g.
\ili{Buru} \it{lakat} ‘long thorns,
particularly on rattan vines or sago palms’ [see \frf{15-03}],
\it{paŋa-t} ‘barb (on spear),
thorn, branch (of antler)’, \it{isinan} ‘sliver,
fine thorn’; \ili{Dela} \it{ŋgaut} ‘thorn’; \ili{Kotos} \ili{Amarasi} \it{aikaʔ}
‘thorn’; \ili{Havu} \it{dudu} ‘thorn’; \ili{Kemak} \it{ai taran} ‘thorn’).
Once again, this is not a replacement innovation,
since reflexes of P(W)MP *tuqəlaŋ ‘candylar bone,
non-fish bone’ are also found throughout the \tsc{Aru} subgroup.
It may be that as we get access to more data in these languages
(for many of which we have only word list data,
and some of that quite limited) we will find reflexes of *tuqəlaŋ
with more specific meanings,
but this remains unknown.